\chapter{Ambiente, Ferramentas e Reprodutibilidade}
\label{chap:ambiente_reprodutibilidade}

\chapterepigraph{``I accept nothing on authority. A hypothesis must be backed by reason, or else it is worthless.''}{Isaac Asimov, \textit{I, Robot} (1950).}
{\small\itshape\color{AccentGray}
A exig\^{e}ncia de Asimov --- que toda hip\'{o}tese seja sustentada por raz\~{a}o, n\~{a}o por autoridade~\cite{asimov1950irobot} --- ressoa diretamente no imperativo freireano de \textit{rigor metodol\'{o}gico}: aceitar premissas sem verifica\c{c}\~{a}o \'{e} uma forma de abdi\c{c}\~{a}o intelectual incompat\'{i}vel com o ato de pensar certo~\cite{freire1996autonomia}. Reprodutibilidade, neste contexto, \'{e} o rigor operacionalizado --- a garantia de que qualquer colaborador pode reconstruir o mesmo ambiente, verificar os mesmos resultados e contestar qualquer afirma\c{c}\~{a}o feita neste livro. O ambiente n\~{a}o \'{e} detalhe operacional; \'{e} condi\c{c}\~{a}o de validade cient\'{i}fica.\par}
\vspace{0.6em}
% ---------------------------------------------------------------
\section{Escopo do cap\'itulo}
% ---------------------------------------------------------------

\begin{quote}\itshape
\textbf{Onde estamos na hist\'oria:} loop interno, fundamentos operacionais.
Antes de construir qualquer agente, o sistema precisa de solo firme:
depend\^encias reproduz\'iveis, credenciais seguras e infraestrutura
declarada como c\'odigo. Este cap\'itulo \'e a base silenciosa sobre a
qual todos os experimentos dos cap\'itulos seguintes se apoiam.
\end{quote}

Este cap\'itulo descreve a configura\c{c}\~ao completa do ambiente de
desenvolvimento adotado no projeto, abrangendo gerenciamento de
depend\^encias Python\index{Python}, credenciais AWS\index{AWS},
provisionamento de infraestrutura
via \textit{Terraform}\index{Terraform|textbf}\index{infraestrutura como c\'odigo|see{Terraform}} e organiza\c{c}\~ao do reposit\'orio. O objetivo
\'e garantir que qualquer colaborador consiga reproduzir os experimentos
descritos ao longo do livro a partir de um ambiente zerado, sem
ambiguidades de vers\~ao ou de configura\c{c}\~ao.

% ---------------------------------------------------------------
\section{Objetivos de aprendizagem}
% ---------------------------------------------------------------

Ao concluir este cap\'itulo, o leitor ser\'a capaz de:

\begin{itemize}[leftmargin=*]
  \item Configurar um ambiente Python~3.11 isolado com todas as
        depend\^encias do projeto fixadas em vers\~ao exata;
  \item Organizar credenciais AWS de forma segura utilizando perfis
        nomeados e vari\'aveis de ambiente;
  \item Provisionar a infraestrutura AWS (Lambda, DynamoDB, SQS,
        OpenSearch, Neptune) via \textit{Terraform} em uma \'unica
        execu\c{c}\~ao reproduz\'ivel;
  \item Executar os \textit{notebooks} Jupyter do projeto e verificar
        os resultados esperados de cada se\c{c}\~ao.
\end{itemize}

% ---------------------------------------------------------------
\section{Conceitos e prerrequisitos}
% ---------------------------------------------------------------

\subsection{Vers\~oes e depend\^encias}

A reprodutibilidade\index{reprodutibilidade|textbf} de experimentos com LLMs \'e comprometida por tr\^es
fontes principais de n\~ao-determinismo: (i)~vers\~ao do modelo de
infer\^encia\index{infer\^encia}, (ii)~vers\~ao das bibliotecas de orquestra\c{c}\~ao e
(iii)~configura\c{c}\~ao do ambiente de execu\c{c}\~ao. Para mitigar esses
riscos, o projeto adota as seguintes pr\'aticas:

\begin{itemize}[leftmargin=*]
  \item Todas as depend\^encias s\~ao declaradas com vers\~ao exata no
        material de instala\c{c}\~ao do projeto, com um comando
        expl\'icito de \texttt{pip install} documentado no
        \texttt{README.md}; para a camada Lambda, as depend\^encias
        continuam separadas em
        \texttt{infraestrutura/modules/lambda\_layer/requirements.txt};
  \item O modelo de infer\^encia \'e referenciado pelo identificador
        versionado \nolinkurl{global.anthropic.claude-sonnet-4-5-20250929-v1:0},
        e n\~ao por alias gen\'erico como \texttt{claude-sonnet-latest};
  \item A vers\~ao do Python (\texttt{3.11.x}) e do \textit{Terraform}
        (\texttt{1.7.x}) s\~ao documentadas no arquivo \texttt{README.md}
        na raiz do reposit\'orio.
\end{itemize}

\subsection{Estrutura do reposit\'orio}

O reposit\'orio \'e organizado segundo o princ\'ipio de separa\c{c}\~ao
de responsabilidades, distinguindo c\'odigo de aplica\c{c}\~ao,
infraestrutura como c\'odigo e artefatos de documenta\c{c}\~ao:

\begin{lstlisting}[language=bash,
  caption={Estrutura de diret\'orios do reposit\'orio.},
  label={lst:repo_structure}]
langchain/
  infraestrutura/          # Terraform (modulos por servico AWS)
    modules/
      lambda/src/          # Codigo das funcoes Lambda
      opensearch/
      neptune/
      dynamodb/
      sqs/
      api_gateway/
  livro/                   # Fonte LaTeX do livro
    chapters/
  decisoes_projeto/        # ADRs (Architecture Decision Records)
  tests/                   # Testes unitarios e de integracao
  langchain.ipynb          # Notebook didatico principal
  demo_api.ipynb           # Notebook de demonstracao da API
  README.md
\end{lstlisting}

% ---------------------------------------------------------------
\section{Implementa\c{c}\~ao orientada por passos}
% ---------------------------------------------------------------

\subsection{Passo 1 --- Clonar o reposit\'orio e criar o ambiente virtual\index{ambiente virtual (Python)|textbf}}

\begin{lstlisting}[language=bash,
  caption={Cria\c{c}\~ao do ambiente virtual Python isolado.},
  label={lst:venv_setup}]
git clone https://github.com/<org>/langchain-rag.git
cd langchain-rag

# Criar e ativar o ambiente virtual
python3.11 -m venv .venv
source .venv/bin/activate          # Linux/macOS
# .venv\Scripts\Activate.ps1       # Windows PowerShell

# Instalar as dependencias do ambiente didatico
pip install --upgrade pip
pip install langgraph langchain langchain-anthropic langchain-community \
            langchain-text-splitters langchain-mcp-adapters faiss-cpu \
            sentence-transformers psycopg2-binary mcp python-dotenv \
            duckduckgo-search wikipedia httpx deepeval
\end{lstlisting}

No estado atual do reposit\'orio, o comando acima \'e a fonte prim\'aria
de instala\c{c}\~ao do ambiente local, conforme o \texttt{README.md}. O
arquivo \texttt{requirements.txt} permanece restrito ao empacotamento da
camada Lambda, e n\~ao ao ambiente did\'atico da raiz do projeto.

\begin{lstlisting}[language=bash,
      caption={Pacotes principais instalados no ambiente local do projeto.},
  label={lst:requirements}]
langgraph
langchain
langchain-anthropic
langchain-community
langchain-text-splitters
langchain-mcp-adapters
faiss-cpu
sentence-transformers
psycopg2-binary
mcp
python-dotenv
duckduckgo-search
wikipedia
httpx
deepeval
\end{lstlisting}

\subsubsection*{Resultado da instala\c{c}\~ao}

Na valida\c{c}\~ao de 2026-05-03, a execu\c{c}\~ao do comando do
C\'odigo~\ref{lst:venv_setup} no Python~3.11.9 n\~ao exigiu novas
instala\c{c}\~oes: os pacotes j\'a estavam presentes no ambiente local.

\begin{lstlisting}[language=bash,
      caption={Sa\'ida resumida da valida\c{c}\~ao das depend\^encias locais.},
      label={lst:cap2_pip_output}]
Requirement already satisfied: langgraph (...)
Requirement already satisfied: langchain (...)
Requirement already satisfied: langchain-anthropic (...)
Requirement already satisfied: langchain-community (...)
Requirement already satisfied: faiss-cpu (...)
Requirement already satisfied: sentence-transformers (...)
Requirement already satisfied: mcp (...)
Requirement already satisfied: deepeval (...)
\end{lstlisting}

\noindent\textbf{Leitura do resultado:} para o leitor, o aspecto
importante n\~ao \'e a lista exaustiva de depend\^encias transitivas, mas
o fato de o comando documentado no projeto resolver o conjunto m\'inimo
de bibliotecas sem erro bloqueante.

\subsection{Passo 2 --- Configurar credenciais AWS}

O acesso aos servi\c{c}os AWS \'e controlado por fun\c{c}\~oes IAM\index{IAM}
(\textit{Identity and Access Management}\index{Identity and Access Management|see{IAM}}). Para desenvolvimento local,
utiliza-se um perfil nomeado com permiss\~oes m\'inimas necess\'arias:

\begin{lstlisting}[language=bash,
  caption={Configura\c{c}\~ao do perfil AWS para desenvolvimento.},
  label={lst:aws_profile}]
aws configure --profile langchain-dev
# AWS Access Key ID:     <ID gerado no IAM>
# AWS Secret Access Key: <CHAVE>
# Default region name:   sa-east-1
# Default output format: json

# Exportar o perfil como variavel de ambiente no shell
export AWS_PROFILE=langchain-dev
\end{lstlisting}

\noindent\textbf{Importante:} credenciais n\~ao devem ser
\textit{hardcoded} no c\'odigo nem versionadas no reposit\'orio.
Recomenda-se incluir \texttt{.env} e \texttt{*.tfvars} no
\texttt{.gitignore}.

\subsubsection*{Resultado da valida\c{c}\~ao da credencial ativa}

Embora o exemplo acima use perfil nomeado, a valida\c{c}\~ao pr\'atica do
ambiente pode ocorrer tamb\'em com credenciais ativas por vari\'aveis de
ambiente ou credenciais padr\~ao da AWS CLI. Em 2026-05-03, o comando
de verifica\c{c}\~ao retornou:

\begin{lstlisting}[language=bash,
      caption={Identidade AWS observada no ambiente de desenvolvimento.},
      label={lst:cap2_aws_identity}]
{
            "UserId": "AIDARU55WWGGKGGKGGNQT",
            "Account": "113677611404",
            "Arn": "arn:aws:iam::113677611404:user/wilian"
}
\end{lstlisting}

\noindent\textbf{Leitura do resultado:} o objetivo do teste n\~ao \'e
for\c{c}ar um mecanismo espec\'ifico de autentica\c{c}\~ao, e sim comprovar
que a identidade ativa pertence \\`a conta esperada antes de qualquer
provisionamento.

\subsection{Passo 3 --- Provisionar a infraestrutura com Terraform}

A infraestrutura AWS \'e descrita como c\'odigo no diret\'orio
\texttt{infraestrutura/}. O provisionamento completo --- incluindo
VPC\index{VPC|textbf}\index{Virtual Private Cloud|see{VPC}},
API Gateway\index{API Gateway|textbf},
filas SQS\index{SQS|textbf}\index{Simple Queue Service|see{SQS}},
fun\c{c}\~oes Lambda\index{Lambda (AWS)|textbf},
tabela DynamoDB\index{DynamoDB|textbf},
dom\'inio OpenSearch\index{OpenSearch|textbf} e cluster Neptune\index{Neptune (AWS)|textbf} ---
\'e executado com:

\begin{lstlisting}[language=bash,
  caption={Provisionamento da infraestrutura via Terraform.},
  label={lst:terraform_apply}]
cd infraestrutura/
terraform init
terraform validate
terraform plan  -out=deploy.tfplan
terraform apply deploy.tfplan
\end{lstlisting}

Ap\'os a conclus\~ao, os \textit{outputs} relevantes --- URL da API,
\textit{endpoints} do OpenSearch e do Neptune, nome da tabela DynamoDB
--- s\~ao exibidos no terminal e podem ser capturados com:

\begin{lstlisting}[language=bash,
  caption={Extra\c{c}\~ao de outputs do Terraform.},
  label={lst:terraform_output}]
terraform output -json > ../model_data.json
\end{lstlisting}

\subsubsection*{Resultado da valida\c{c}\~ao do Terraform}

A execu\c{c}\~ao observada em 2026-05-03 confirmou o fluxo correto:
\texttt{terraform validate} s\'o \'e conclusivo ap\'os
\texttt{terraform init}, pois os m\'odulos precisam estar instalados
localmente.

\begin{lstlisting}[language=bash,
      caption={Sa\'ida resumida do \\texttt{terraform init} local.},
      label={lst:cap2_terraform_init_output}]
Initializing modules...
- api_gateway in modules\api_gateway
- cloudfront_frontend in modules\cloudfront_frontend
- dynamodb in modules\dynamodb
...
Terraform has been successfully initialized!
\end{lstlisting}

\begin{lstlisting}[language=bash,
      caption={Sa\'ida do \\texttt{terraform validate} ap\'os a inicializa\c{c}\~ao.},
      label={lst:cap2_terraform_validate_output}]
Success! The configuration is valid.
\end{lstlisting}

\noindent\textbf{Leitura do resultado:} esse detalhe operacional \'e
importante para o leitor porque evita um falso diagn\'ostico. Um erro de
\textit{module not installed} nessa etapa n\~ao indica falha da
infraestrutura, apenas aus\^encia da inicializa\c{c}\~ao pr\'evia.

\subsection{Passo 4 --- Executar os notebooks}

Com o ambiente configurado e a infraestrutura provisionada, os
\textit{notebooks} podem ser executados diretamente no VS Code com a
extens\~ao \textit{Jupyter}, que foi o fluxo validado nesta revis\~ao. Se
o leitor preferir um servidor local no navegador, o JupyterLab deve ser
instalado explicitamente no ambiente Python:

\begin{lstlisting}[language=bash,
      caption={Fluxo opcional para iniciar um servidor local do JupyterLab.},
  label={lst:jupyter_start}]
python -m pip install jupyterlab
jupyter lab --no-browser --port=8888
\end{lstlisting}

Recomenda-se executar os notebooks na ordem natural dos cap\'itulos ---
\texttt{langchain.ipynb} (conceitos) e \texttt{demo\_api.ipynb}
(demonstra\c{c}\~ao da API) --- pois as c\'elulas iniciais de cada
\textit{notebook} definem constantes e fun\c{c}\~oes utilizadas
nas c\'elulas subsequentes.

\subsubsection*{Resultado da execu\c{c}\~ao no VS Code}

No ambiente validado em 2026-05-03, a primeira c\'elula de c\'odigo do
arquivo \texttt{langchain.ipynb} executou com sucesso no kernel
Python~3.11.9 do VS Code e retornou:

\begin{lstlisting}[language=bash,
      caption={Sa\'ida da primeira c\'elula do notebook principal no VS Code.},
      label={lst:cap2_notebook_output}]
Dependencias instaladas com sucesso!
\end{lstlisting}

\noindent\textbf{Leitura do resultado:} para o leitor, essa evid\^encia
mostra que o fluxo principal do projeto hoje \'e a execu\c{c}\~ao do
notebook no editor, sem depend\^encia obrigat\'oria de um processo
externo de JupyterLab.

% ---------------------------------------------------------------
\section{Plano de testes do cap\'itulo}
% ---------------------------------------------------------------

Esta se\c{c}\~ao apresenta os comandos de valida\c{c}\~ao efetivamente
executados no ambiente do projeto. Sempre que poss\'ivel, o c\'odigo do
teste aparece junto da sa\'ida observada e de um coment\'ario curto para
orientar a interpreta\c{c}\~ao do leitor.

\subsection{Teste funcional}

\begin{itemize}[leftmargin=*]
  \item \textbf{TC-01:} Executar o comando de instala\c{c}\~ao
        documentado no \texttt{README.md} e verificar aus\^encia de
        erros bloqueantes no \textit{resolver} do \texttt{pip}.
  \item \textbf{TC-02:} Executar \texttt{python -c "import boto3,
        langchain, langgraph, faiss"} e verificar que nenhum
        \texttt{ImportError} \'e lan\c{c}ado.
  \item \textbf{TC-03:} Executar \texttt{aws sts get-caller-identity}
        com a credencial ativa (perfil nomeado ou ambiente padr\~ao) e
        verificar que o ARN retornado pertence \`a conta correta.
  \item \textbf{TC-04:} Executar \texttt{terraform init} seguido de
        \texttt{terraform validate} no diret\'orio de infraestrutura e
        verificar mensagem
        \texttt{Success! The configuration is valid.}
\end{itemize}

\subsubsection*{Evid\^encias de execu\c{c}\~ao (2026-05-03)}

\paragraph{TC-01 --- instala\c{c}\~ao do ambiente}

O teste reutiliza exatamente o comando documentado no
C\'odigo~\ref{lst:venv_setup}. A sa\'ida resumida foi apresentada no
C\'odigo~\ref{lst:cap2_pip_output}. O ponto importante para o leitor \'e
que a instala\c{c}\~ao n\~ao falhou por conflito de depend\^encias.

\paragraph{TC-02 --- importa\c{c}\~oes essenciais}

\begin{lstlisting}[language=Python,
  caption={TC-02 --- verifica\c{c}\~ao m\'inima de importa\c{c}\~oes do ambiente.},
  label={lst:cap2_import_check}]
import boto3, langchain, langgraph, faiss
print("IMPORTS_OK")
\end{lstlisting}

\begin{lstlisting}[language=bash,
  caption={TC-02 --- sa\'ida da verifica\c{c}\~ao de importa\c{c}\~oes.},
  label={lst:cap2_import_output}]
IMPORTS_OK
\end{lstlisting}

\noindent\textbf{Leitura do resultado:} esse teste n\~ao prova que todo
o projeto est\'a funcional, mas confirma que o n\'ucleo do ambiente
Python foi carregado sem erro de m\'odulo ausente.

\paragraph{TC-03 --- identidade AWS ativa}

O comando executado foi \texttt{aws sts get-caller-identity}; a sa\'ida
observada consta no C\'odigo~\ref{lst:cap2_aws_identity}. O leitor deve
verificar especialmente o campo \texttt{Arn}, pois ele identifica com
clareza a entidade autenticada.

\paragraph{TC-04 --- valida\c{c}\~ao do Terraform}

O teste foi decomposto em duas etapas, conforme os
C\'odigos~\ref{lst:cap2_terraform_init_output} e
\ref{lst:cap2_terraform_validate_output}. Essa decomposi\c{c}\~ao evita
confundir falha de inicializa\c{c}\~ao com erro sem\^antico de
configura\c{c}\~ao.

\subsection{Teste de desempenho}

\begin{itemize}[leftmargin=*]
  \item \textbf{TP-01:} Medir o tempo total de \texttt{terraform apply}
        em ambiente limpo; crit\'erio: conclus\~ao em menos de
        20\,minutos para o conjunto completo de m\'odulos.
  \item \textbf{TP-02:} Medir o tempo de inicializa\c{c}\~ao do
      kernel do notebook no VS Code at\'e a execu\c{c}\~ao da primeira
      c\'elula de importa\c{c}\~ao; crit\'erio: menos de 30\,s.
\end{itemize}

\subsubsection*{Observa\c{c}\~ao sobre TP-02}

Na valida\c{c}\~ao de 2026-05-03, o notebook principal executou a
primeira c\'elula com sucesso no VS Code. O ambiente n\~ao possu\'ia
\texttt{jupyter-lab} pr\'e-instalado, raz\~ao pela qual o crit\'erio foi
explicitado em termos do fluxo real de uso do projeto, e n\~ao de um
servidor local opcional.

\subsection{Teste de qualidade}

\begin{itemize}[leftmargin=*]
  \item \textbf{TQ-01:} Verificar que o arquivo
        \texttt{.gitignore} exclui \texttt{*.tfstate},
        \texttt{*.tfvars} e \texttt{.env} do versionamento.
  \item \textbf{TQ-02:} Confirmar que nenhuma chave de acesso AWS
                        com padr\~ao realista \texttt{AKIA[0-9A-Z]\{16\}} est\'a
                        presente na \textit{working tree} rastreada pelo Git.
\end{itemize}

\subsubsection*{Evid\^encias de qualidade}

\begin{lstlisting}[language=bash,
      caption={TQ-01 --- fragmento relevante do \texttt{.gitignore}.},
      label={lst:cap2_gitignore_output}]
.env
**/.terraform/
infraestrutura/*.tfstate
infraestrutura/*.tfstate.backup
infraestrutura/secrets.auto.tfvars
\end{lstlisting}

\begin{lstlisting}[language=bash,
      caption={TQ-02 --- busca por padr\~ao de chave AWS na \textit{working tree}.},
      label={lst:cap2_akia_output}]
PADRAO_NAO_ENCONTRADO
\end{lstlisting}

\noindent\textbf{Leitura do resultado:} a primeira tentativa de busca
textual simples por \texttt{AKIA} produziu falso positivo no pr\'oprio
texto do livro. Por isso, o teste foi refinado para um padr\~ao de chave
mais preciso, reduzindo ru\'ido e tornando a checagem reprodut\'ivel.

% ---------------------------------------------------------------
\section{Crit\'erios de aceite}
% ---------------------------------------------------------------

\begin{itemize}[leftmargin=*]
  \item TC-01 a TC-04 conclu\'idos sem erros em um ambiente
        Windows~11 com Python~3.11 e Terraform~1.7.
  \item TP-01: \texttt{terraform apply} conclu\'ido em menos de
        20\,minutos (evidenciar com \textit{log} de execu\c{c}\~ao).
  \item TQ-01 e TQ-02 aprovados; nenhum padr\~ao plaus\'ivel de
        credencial AWS detectado na \textit{working tree}.
\end{itemize}

% ---------------------------------------------------------------
\section{Espa\c{c}o para expans\~ao iterativa}
% ---------------------------------------------------------------

\begin{itemize}[leftmargin=*]
  \item \textbf{v2.1} --- Adicionar \textit{devcontainer} Docker para
        eliminar depend\^encia de instala\c{c}\~ao local de Python e
        Terraform.
  \item \textbf{v2.2} --- Incluir pipeline CI/CD (GitHub Actions) que
        execute \texttt{terraform validate} e os testes unit\'arios a
        cada \textit{pull request}.
  \item \textbf{Lacuna} --- Avaliar o uso de \textit{Terraform Workspaces}
        para isolar ambientes de desenvolvimento, homologa\c{c}\~ao e
        produ\c{c}\~ao sem duplicar c\'odigo de infraestrutura.
\end{itemize}
